\section{Study Design and Data}

The study focuses on Kazakhstan as a data-scarce urban setting in which official city totals are available but true grid-level population labels are not. This makes the national setting methodologically central rather than incidental. The framework is designed for a context where city-level accounting can be anchored to verified totals, while intra-urban structure must be inferred from open geospatial evidence under explicit uncertainty and limitation boundaries.

The broader deterministic baseline is organized around a 10-city reference registry with official totals available for calibrated inference: Almaty, Astana, Shymkent, Semey, Kurchatov, Taraz, Uralsk, Petropavlovsk, Ust Kamenogorsk, and Ridder. Within that registry, 8 cities form the validated baseline batch with the full combination of official totals, feature generation, and baseline evidence support: Almaty, Astana, Shymkent, Semey, Kurchatov, Taraz, Uralsk, and Petropavlovsk. Kurchatov and Ridder remain important for calibrated runtime support or official-total registry coverage, even though they do not carry the full deeper validation stack. Semey is additionally retained as a smoke-passed example city in the reference package. This is not an inconsistency. It is a layered evidence design in which the registry, the baseline validation batch, and the uncertainty-evidence subset answer different scientific questions at different depths.

Official city totals act as calibration anchors throughout the framework. The canonical reference is \texttt{data/external/city\_population\_reference\_v2.csv}, with supported totals fixed to the reference date 2026-02-01. In the reference package, Almaty, Astana, and Shymkent are tied to Bureau of National Statistics regional statistics page values, while the other supported cities are compiled through cached BNS regional monthly workbooks. The companion registry \texttt{data/external/city\_status\_registry\_v2.csv} records which cities belong to the validated baseline, the smoke-tested examples, and the broader support frame.

The feature layer is derived from OpenStreetMap and grouped into built form, floor-area proxies, transport structure, and point-of-interest or service-access variables. This design reflects the bounded premise of the framework: in the absence of true cell-level census labels, the target is not truth reconstruction but a transparent calibrated proxy surface supported by reproducible open-data inputs. The production grid is fixed to 500 m for both analytical stages.

The study design also distinguishes between the broader baseline registry and the narrower reliability-evidence subset. District benchmark comparison is available for Almaty, Astana, and Shymkent. Qualitative plausibility review is locked to Almaty and Astana as two deliberately contrasting visual reading contexts. External structural comparison is conducted against WorldPop and GHS-POP. OSM completeness is retained as input-support context rather than as a performance label.

The uncertainty-aware extension then focuses on a four-city inference package: Almaty, Astana, Semey, and Shymkent. This subset is used to test reliability-aware interpretation, confidence-band separation, interval behavior, and hotspot screening. It does not replace the broader deterministic 10-city reference frame. Instead, the relationship is explicit: the deterministic baseline uses the broader registry and validation stack, whereas the uncertainty-aware extension uses the four-city package to assess where the calibrated proxy surface is stronger, weaker, or caution-heavier.

Seen together, the study design therefore behaves like a staged evidence ladder. The widest registry supports reproducibility and coverage. The eight-city baseline supports the main deterministic claims. The three district benchmark cities and two qualitative plausibility cities support specific kinds of credibility checks. The four-city uncertainty package supports reliability-aware interpretation. Each subset is narrower for a reason, and each reason is traceable to the claim boundary it is supposed to support.

\input{tables_tex/table2_study_coverage_summary.tex}
